\section{A local Cauchy theorem on arbitrarily shrinking intervals}
\label{sec:local-cauchy}

The scalar part of the preceding proof admits parameters that change with
$T$.  This produces a local limit theorem rather than only a weak limit.
It is also the most direct statement of the result for a single classical
special function.

Fix one of the completed Dirichlet functions $F$ from
Section~\ref{sec:setup}, and put $g=F'/F$.  On the real line away from
the zeros of $F$, the function $g$ is real.  Let $u<v$ be fixed real
numbers and let $U_T$ be uniform on $[uT,vT]$.  Define
\begin{equation}\label{eq:normalized-random-variable}
 Z_T=\frac{2}{\log T}\frac{F'(U_T)}{F(U_T)}.
\end{equation}
The undefined values at zeros are a null event and may be assigned
arbitrarily.

\begin{theorem}[Arbitrarily fine local Cauchy law]\label{thm:local-cauchy}
Let $\alpha_T<\beta_T$ be any two real-valued functions of $T$ such that
\[
 \alpha_T\longrightarrow c,\qquad\beta_T\longrightarrow c
       \quad\text{for some fixed }c\in\mathbb R.
\]
There is no lower bound on the positive width $\beta_T-\alpha_T$.
Then
\begin{equation}\label{eq:local-cauchy-law}
 \boxed{\displaystyle
  \Pr\{\alpha_T<Z_T<\beta_T\}
      \sim\frac{\beta_T-\alpha_T}{\pi(1+c^2)}.}
\end{equation}
Equivalently, if $a_T<b_T$, $\delta_T=b_T-a_T>0$, and
$a_T/\log T,b_T/\log T\to d\in\mathbb R$, then
\begin{equation}\label{eq:raw-local-law}
 \boxed{\displaystyle
 \left|\left\{t\in[uT,vT]:a_T<g(t)<b_T\right\}\right|
 \sim\frac{2(v-u)}{\pi(1+4d^2)}
                    \frac{\delta_T T}{\log T}.}
\end{equation}
Vertical bars denote Lebesgue measure, not a number of sample points.
All characters and conductors remain fixed.
\end{theorem}

\begin{proof}
We prove \eqref{eq:raw-local-law}.  For each $T$, define in the upper
half-plane
\begin{equation}\label{eq:moving-band-log}
 f_T(z)=\int_{a_T}^{b_T}\frac{ds}{g(z)-s}.
\end{equation}
Since $\operatorname{Im}g(z)<0$, the integrand has positive imaginary
part.  This is the analytic branch of
$\log[(g(z)-a_T)/(g(z)-b_T)]$ with imaginary part between $0$ and $\pi$.
For almost every real $x$, its limiting imaginary part is
\[
 \operatorname{Im}f_T(x+i0)
        =\pi\,\boldsymbol1_{\{a_T<g(x)<b_T\}}.
\]
For example this follows directly by taking the limit of
\[
 \operatorname{Im}\frac1{g(z)-s}
        =\frac{-\operatorname{Im}g(z)}
                     {(\operatorname{Re}g(z)-s)^2+(\operatorname{Im}g(z))^2}
\]
in the finite $s$-integral.  The exceptional endpoints and poles are
discrete.  The bounded harmonic Poisson formula therefore applies
exactly, for every $T$.

By Lemma~\ref{ut:radial}, uniformly on compact subsets of the upper
half-plane,
\[
 \frac{g(Tz)}{\log T}\longrightarrow-\frac i2.
\]
Every $s$ between $a_T$ and $b_T$ satisfies $s/\log T\to d$, uniformly
over that interval.  Dividing \eqref{eq:moving-band-log} by its length
thus gives
\begin{equation}\label{eq:band-harmonic-limit}
 \frac{\log T}{\delta_T}f_T(Tz)
       =\frac1{\delta_T}\int_{a_T}^{b_T}
               \frac{ds}{g(Tz)/\log T-s/\log T}
       \longrightarrow\frac1{-d-i/2}.
\end{equation}
The denominators stay uniformly away from zero.  This averaging argument
does not divide an uncontrolled error by $\delta_T$: the error is uniform
\emph{inside the integral}, and its integral is $\delta_T$ times that
uniform error.  Hence no restriction on how quickly $\delta_T$ tends to
zero enters the proof.

Introduce the positive measures
\[
 d\mu_T(x)=\frac{\log T}{\delta_T}
       \boldsymbol1_{\{a_T<g(Tx)<b_T\}}\,dx.
\]
Their Poisson transforms, by \eqref{eq:band-harmonic-limit}, tend at every
upper-half-plane point to
\[
 \operatorname{Im}\frac1{-d-i/2}=\frac2{1+4d^2}.
\]
The compactness, tail bound, and Poisson uniqueness argument in the
proof of Theorem~\ref{ut:main} uses only positivity and convergence of
these transforms.  It applies to the present $T$-dependent measures as
well, giving
\[
 d\mu_T(x)\longrightarrow\frac2{\pi(1+4d^2)}\,dx
\]
vaguely.  In more detail, the transform at $iR$ bounds the mass on
$[-R,R]$; the kernel inequality in that proof bounds all tail transforms
by $O(1/R)$; every subsequential limit then has the displayed constant
Poisson transform, determining it uniquely.  The continuous limiting
density permits taking the mass of $[u,v]$.  This is
\eqref{eq:raw-local-law}.

Finally set $a_T=\alpha_T\log T/2$ and
$b_T=\beta_T\log T/2$.  Then $d=c/2$.  Divide the resulting measure
by the length $(v-u)T$ of the sampling interval to obtain
\eqref{eq:local-cauchy-law}.
\end{proof}

\begin{corollary}[Exact asymptotic for rare near-stationarity]
\label{cor:rare-flatness}
For every function $\varepsilon(T)>0$ with
$\varepsilon(T)=o(\log T)$,
\begin{equation}\label{eq:near-stationarity}
 \boxed{\displaystyle
 \left|\left\{t\in[T,2T]:
           \left|\frac{F'(t)}{F(t)}\right|<\varepsilon(T)\right\}\right|
        \sim\frac{4\varepsilon(T)T}{\pi\log T}.}
\end{equation}
This applies, for example, to $\varepsilon(T)=e^{-T}$.
For the Riemann input, $F(t)=\Xi(t)=\xi(1/2+it)$.
\end{corollary}
\begin{proof}
Take $a_T=-\varepsilon(T)$ and $b_T=\varepsilon(T)$, so $d=0$ and
$\delta_T=2\varepsilon(T)$, in \eqref{eq:raw-local-law}.
\end{proof}

\paragraph{What is stronger than a weak Cauchy law.}
Weak convergence controls probabilities of fixed intervals whose
boundaries have limiting measure zero.  It supplies no relative error
for an interval with length tending to zero, especially a length as
small as $e^{-T}$.  Theorem~\ref{thm:local-cauchy} gives that relative
asymptotic directly.  It concerns a continuous sampling variable; it
does not claim that a discrete sampling grid resolves exponentially
short intervals.

Scalar weak Cauchy limits for zeta logarithmic derivatives have already
been studied by Sodin and, more recently, Najnudel and
Nikeghbali~\cite{sodin,najnudel}.  The latter's Theorem~1.1 gives the
weak Cauchy law, while its introduction distinguishes the available
comparison estimate from a rate of convergence to that law.  Our
proof uses the established real-zero setting and a positive bounded
phase to obtain the stronger moving-interval statement.  We claim
the derivation and comparison recorded here, not the invention of
the underlying Cauchy phenomenon.
